El backend se despliega en local mediante \texttt{docker-compose},
conforme al \cref{rnf:deployment}. Esta sección documenta las decisiones
técnicas concretas que materializan ese despliegue; el manual operativo
paso a paso se incluye en \cref{AP:INSTALLATION}.

\begin{itemize}

  \item \textbf{Imagen Docker compartida para los servicios Python.}
  \acs{api}, \textit{worker} de Celery y Celery \textit{beat} ejecutan
  código de la misma aplicación Django y comparten una única imagen
  construida a partir de \texttt{python:3.12-slim}, con las dependencias
  de sistema necesarias para los flujos críticos
  (\texttt{libzbar0} para \acs{qr}, \texttt{poppler-utils} para
  rasterización de \acsp{pdf}, \texttt{tesseract-ocr-spa} para \acs{ocr},
  \texttt{libpq-dev} para psycopg). El contenedor se ejecuta con un
  usuario sin privilegios de \textit{root} y el \texttt{Dockerfile}
  instala dependencias antes que el código para optimizar el caché de
  capas.

  \item \textbf{Configuración por variables de entorno.} Toda la
  configuración específica del entorno (credenciales de base de datos,
  claves de cifrado, parámetros operativos, \acs{uri} del servidor
  \acs{smtp}) se externaliza a un fichero \texttt{.env} que cada
  despliegue puede personalizar sin reconstruir la imagen.

  \item \textbf{Persistencia mediante volúmenes Docker.} PostgreSQL,
  Redis, MinIO y la bandeja de ingesta usan volúmenes Docker nombrados
  que sobreviven a los reinicios de los contenedores.

  \item \textbf{Arranque ordenado mediante \textit{healthchecks}.} Cada
  servicio de infraestructura (PostgreSQL, Redis, MinIO) define un
  \textit{healthcheck} propio; los servicios de aplicación declaran
  dependencias con condición \texttt{service\_healthy}, lo que garantiza
  que el \acs{api} y los \textit{workers} sólo arrancan cuando sus
  dependencias están operativas.

  \item \textbf{Perfiles para separar entornos.} Los servicios de
  monitorización descritos en \cref{SS:DES-COMPONENTS} viven en un
  perfil aparte (\texttt{monitoring}), no se levantan por defecto y se
  activan bajo demanda con
  \texttt{docker-compose --profile monitoring up}.

\end{itemize}
